\section{Authority: is this within Congress's power and consistent with existing law?}
\label{sec:authority}

The second question is whether Congress may act here at all, and whether the bill
respects the law already on the books. A clean answer disarms both the
constitutional objection and the federalism objection before they are raised.

Congress's authority is the same it has exercised over medical devices for decades.
H. R. 9510 amends the Federal Food, Drug, and Cosmetic Act, 21 U.S.C. 301 et seq.,
an exercise of the commerce power, and adds a single new section 515D (21 U.S.C.
360e-5) together with ten conforming amendments that thread the new requirement
through the existing device provisions, from the definition of a device in section
321(h) to the State preemption rule in section 360k \cite{fdcact,
Kawchak2026HR9510Bill}. The bill does not invent a new regulatory regime; it places a
sequencing requirement inside the one Congress already built and the FDA already
administers.

The bill is careful with existing law. Its rule of construction states that nothing
in section 515D displaces the investigational-device exemption of section 520(g),
the human-subjects protections of 21 CFR part 50, or the investigational new drug
framework of 21 CFR part 312, and that nothing reclassifies any device. Those are the
same frameworks the National Platform adapts rather than replaces
\cite{Kawchak2026Adaption21CFRPart, Kawchak2026Adaption21CFRPartb,
Kawchak2026AdaptionICHHarmonisedGuideline}. The closest existing analogues confirm
that the bill fills a gap rather than overlapping: section 515C predetermined change
control plans and the FDA's guidance on AI-enabled device software are the nearest
precedents, but both are voluntary and neither requires verification before
generation, while the strongest Federal precedent for human oversight of an
algorithm, the Medicare Advantage rule on the use of algorithms in coverage
determinations, governs payment rather than a robot at the bedside
\cite{cms2024aiguardrails}.

Federalism is handled by a savings clause, not by silence. The new section 360k(c)
preserves State requirements that a licensed human retain authority to monitor,
override, or terminate an AI or robotic system, and that audit records be kept, so
long as those requirements are not different from or in addition to the Federal
floor. \autoref{fig:authority} shows authority flowing from the commerce power
through the FD\&C Act to a Federal floor that leaves State human-over-AI law standing.

\begin{psgfig}
\node[psgone] (cc) at (0,0) {Commerce in\\medical devices};
\node[psgevid] (fdc) at (4.0,0) {FD\&C Act\\301 et seq.};
\node[psgact] (fl) at (8.0,0) {Section 515D\\Federal floor};
\node[psgthree] (sv) at (8.0,-2.0) {360k(c) savings:\\State human-over-AI};
\draw[psgedge] (cc)--(fdc); \draw[psgedge] (fdc)--(fl);
\draw[psgedge] (fl)--(sv);
\end{psgfig}
\psgcaption{\textbf{Figure 3.} Authority flows from the commerce power through the
FD\&C Act to a Federal floor in section 515D that preserves State human-over-AI law;
the colored TikZ reproduction of catalog diagram 09.}

\needspace{9\baselineskip}\begin{table}[H]
\footnotesize
\begin{tabularx}{\textwidth}{@{}L{4.2cm} Y@{}}
\toprule
\textbf{Rule-of-construction clause} & \textbf{The existing law it preserves}\\
\midrule
Investigational-device exemption & Section 520(g) of the FD\&C Act is undisturbed.\\
Human-subjects protection & 21 CFR part 50 continues to govern consent and protection.\\
Investigational new drug framework & 21 CFR part 312 continues to govern the trial.\\
Change control outside Physical AI & Section 515C authority is unchanged for other devices.\\
State human-over-AI requirements & New section 360k(c) preserves them as a floor, not a ceiling.\\
\bottomrule
\end{tabularx}
\caption{Authority with restraint: each rule-of-construction clause names the
existing law the bill leaves in place.}
\label{tab:authority}
\end{table}

Authority is answered by precedent and restraint: a familiar commerce-power amendment
to a statute the FDA already runs, with an explicit promise about what it does not
touch.
